% year: תשע"ט
% semester: א
% moed: ב
% number: 6א
% points: 10
% categories: ivt, func-analysis
% source: exams/מבחנים/תשעט/מועד ב תשעט.pdf

\begin{question}
הוכיחו כי למשוואה $e^{2x}-x-5=0$ יש בדיוק שני פתרונות ממשיים.
\end{question}

\begin{hint}
בדקו את המונוטוניות של $g(x)=e^{2x}-x-5$ בעזרת הנגזרת, ואת ערכה בנקודת המינימום. קיום – לפי משפט ערך הביניים, יחידות בכל קטע – לפי מונוטוניות.
\end{hint}

\begin{solution}
נגדיר $g(x)=e^{2x}-x-5$, רציפה וגזירה בכל $\R$, עם $g'(x)=2e^{2x}-1$. $g'(x)=0\iff e^{2x}=\frac12\iff x_0=-\frac{\ln2}{2}$. עבור $x<x_0$: $g'<0$, $g$ יורדת ממש; עבור $x>x_0$: $g'>0$, $g$ עולה ממש. ערך המינימום:
\[ g(x_0)=\frac12+\frac{\ln2}{2}-5\approx-4.15<0. \]
\textbf{קיום:} $\lim_{x\to-\infty}g(x)=+\infty$ (כי $e^{2x}\to0$ ו-$-x\to+\infty$), ולמשל $g(-6)=e^{-12}+1>0$. כמו כן $g(2)=e^4-7>0$. לפי משפט ערך הביניים יש שורש ב-$(-6,x_0)$ ושורש ב-$(x_0,2)$.

\textbf{יחידות:} בכל אחד מהקטעים $(-\infty,x_0]$ ו-$[x_0,\infty)$ הפונקציה מונוטונית ממש ולכן חח"ע, כך שיש בכל אחד לכל היותר שורש אחד (ו-$x_0$ עצמו אינו שורש). לכן למשוואה בדיוק שני פתרונות ממשיים.
\end{solution}
